% The long abstract, used by the ACM version (preplang.tex).
Batch data-preparation pipelines are compiled once and run again on every new batch: next
month's file has the same layout and different values. A compiler that reads the input data at
compile time can prove a pipeline's contracts and delete the checks and operations the data makes
useless, but its proofs are about one file. We present \PL, a compiler for table and image
preprocessing pipelines that computes an abstract dataset state from the data, or from its
metadata alone, decides every contract as proven, violated or deferred, and optimizes, together
with three ways to carry its proofs to new batches. \emph{Data certificates}, found by re-analysing
the program under counterfactually weakened input facts, are sets of input facts that the proofs
and removals depend on, 1-minimal for programs without derived columns. \emph{Selective deoptimization} keeps every removed check and
operation in the generated program with its own item certificate: a new batch runs the optimized
program, and only the removals whose certificates it breaks run again; a batch is refused only for
facts no run-time code can stand in for, such as a column's type. We prove this exact for every
admitted batch and mechanize the proof in Lean~4, end to end for a core of the language. \emph{Composition} lets a compiled stage
publish a guarantee for each table it writes; checking it against the next stage's certificates
shows which downstream checks can never fire and writes the upstream contracts that would close
the gap. On \RealRuns{} real monthly batches of three public datasets, selective deoptimization
handles \RealDeoptCorrect{} batches exactly as a fresh compilation would and none
unsoundly, where Great Expectations and pandera, set up on the same data, refuse
\RealValidatorStopMin--\RealValidatorStopMax\% of the months a fresh compilation runs, and TFDV,
which checks no numeric ranges, admits months on which the optimized program does not do what a fresh
compilation does. Pandas scripts from GitHub, lifted into \PL, reproduce their own
results on new batches: \CorpusBothSame{} of \CorpusBothRuns{} runs of \CorpusBothLifted{} held-out
scripts for two Kaggle tasks, and all \DriftRuns{} monthly runs of \DriftCompiled{} scripts written by
others for the traffic and bike data; none differ. Because selective deoptimization never relies on
the compile-time facts being true, \PL can also compile against data the analyst may not see. We prove,
in Lean for a core, that for programs without derived columns or contracts that read values, what the
analysis decides depends on the data only through which cells, cut at the program's own numbers, hold a
row; the data owner releases a differentially private histogram over
exactly these cells, per row or per person, and everything the compiler writes is
$\varepsilon$-differentially private. \PrivRuns{} runs of \PrivCompilations{} private compilations, from noisy,
absent or random facts, were all exact.
